\jans{1}{\ansul{\(- \frac{3}{2}\)}}
【分析】根据题意\(\overrightarrow{m}\bot\overrightarrow{n}\)得到\(\left( \overrightarrow{a} + \overrightarrow{b} \right) \cdot \allowbreak\left( \overrightarrow{a} + \lambda\overrightarrow{b} \right) = 0\)，从而可求出\(\lambda\)的值.
【详解】由\(\overrightarrow{m}\bot\overrightarrow{n}\)，得\(\left( \overrightarrow{a} + \overrightarrow{b} \right) \cdot \allowbreak\left( \overrightarrow{a} + \lambda\overrightarrow{b} \right) = 0\)，即\({\overrightarrow{a}}^{2} + (\lambda + 1) \cdot \overrightarrow{a} \cdot \overrightarrow{b} + \lambda{\overrightarrow{b}}^{2} = 0\)，即
\(\left| \overrightarrow{a} \right|^{2} + (\lambda + 1) \cdot \allowbreak\left| \overrightarrow{a} \right| \cdot \allowbreak\left| \overrightarrow{b} \right| \cdot \cos\left\langle \overrightarrow{a},\overrightarrow{b} \right\rangle + \lambda\left| \overrightarrow{b} \right|^{2} = 0\)，所以\(18 + (\lambda + 1) \times 3\sqrt{2} \times 4 \times \cos 135{^\circ} + 16\lambda = 0\)，即4\emph{λ}＋6＝0，∴\emph{λ}＝\(- \frac{3}{2}\)．故答案为：\(- \frac{3}{2}\).
\jans{2}{\ansul{C}}
【分析】利用向量加法的三角形法则先求前两项的和，再由减法的几何意义作差即可．
【详解】解：\(\overrightarrow{AB} + \overrightarrow{BC} = \overrightarrow{AC}\)，故\(\overrightarrow{AB} + \overrightarrow{BC} - \overrightarrow{AD} = \overrightarrow{AC} - \overrightarrow{AD} = \overrightarrow{DC}\)．故选：C．
\jans{3}{\ansul{\(\overrightarrow{b} - \overrightarrow{c}\)；\(\overrightarrow{a} + \overrightarrow{b} + \overrightarrow{c}\)}}
【分析】利用空间向量加减法的几何意义，把目标向量沿棱向量作和差分解．
【详解】解：\(\overrightarrow{A_{1}D} = \overrightarrow{AD} - \overrightarrow{AA_{1}} = \overrightarrow{b} - \overrightarrow{c}\)；\(\overrightarrow{AC_{1}} = \overrightarrow{AB} + \overrightarrow{BC} + \overrightarrow{CC_{1}} = \overrightarrow{a} + \overrightarrow{b} + \overrightarrow{c}\)．
\jans{4}{\ansul{A}}
【分析】根据零向量、相反向量、共线向量与相等向量的定义逐一判断即可．
【详解】解：由规定，零向量与任意向量平行（共线），①对；由相反向量的定义，②对；共线向量所在的直线互相平行或重合，不必重合，③错；相等向量只需方向相同且模相等，与起点无关，④错．故选：A．
\jans{5}{\ansul{B}}
【分析】利用数量积与夹角的关系，注意反向共线时夹角为180°的特例，判断充分性与必要性．
【详解】解：若\(\overrightarrow{a} \cdot \overrightarrow{b} < 0\)，则\(\cos\langle \overrightarrow{a},\overrightarrow{b} \rangle < 0\)，\(\langle \overrightarrow{a},\overrightarrow{b} \rangle \in \left( 90^{\circ},180^{\circ} \right]\)，当两向量反向共线时夹角为\(180^{\circ}\)，不是钝角，充分性不成立；反之，夹角为钝角时\(\cos\langle \overrightarrow{a},\overrightarrow{b} \rangle < 0\)，必有\(\overrightarrow{a} \cdot \overrightarrow{b} < 0\)，必要性成立．故选：B．
\jans{6}{\ansul{\(120^{\circ}\)}}
【分析】把\(\overrightarrow{A_{1}B}\)与\(\overrightarrow{BC_{1}}\)用在同一顶点处的三条棱向量表示，用数量积求夹角．
【详解】解：不妨设棱长为1，\(\overrightarrow{AB} = \overrightarrow{a}\)，\(\overrightarrow{AD} = \overrightarrow{b}\)，\(\overrightarrow{AA_{1}} = \overrightarrow{c}\)，则\(\overrightarrow{a}\)，\(\overrightarrow{b}\)，\(\overrightarrow{c}\)两两垂直且模均为1．于是\(\overrightarrow{A_{1}B} = \overrightarrow{a} - \overrightarrow{c}\)，\(\overrightarrow{BC_{1}} = \overrightarrow{b} + \overrightarrow{c}\)，\(\overrightarrow{A_{1}B} \cdot \overrightarrow{BC_{1}} = \allowbreak\left( \overrightarrow{a} - \overrightarrow{c} \right) \cdot \allowbreak\left( \overrightarrow{b} + \overrightarrow{c} \right) = \overrightarrow{a} \cdot \overrightarrow{b} + \overrightarrow{a} \cdot \overrightarrow{c} - \overrightarrow{c} \cdot \overrightarrow{b} - {\overrightarrow{c}}^{2} = - 1\)，又\(\left| \overrightarrow{A_{1}B} \right| = \sqrt{2}\)，\(\left| \overrightarrow{BC_{1}} \right| = \sqrt{2}\)，所以\(\cos\langle \overrightarrow{A_{1}B},\overrightarrow{BC_{1}} \rangle = \frac{- 1}{\sqrt{2} \times \sqrt{2}} = - \frac{1}{2}\)．因为\(\langle \overrightarrow{A_{1}B},\overrightarrow{BC_{1}} \rangle \in \lbrack 0,\pi\rbrack\)，所以\(\overrightarrow{A_{1}B}\)与\(\overrightarrow{BC_{1}}\)的夹角为\(120^{\circ}\)．
\jans{7}{\ansul{\(- \frac{3}{2}\)；\(- \frac{3}{4}\overrightarrow{b}\)}}
【分析】由投影向量（条目8）的定义直接计算：投影向量＝模长×夹角余弦×单位向量．
【详解】解：\(\left| \overrightarrow{a} \right|\cos\langle \overrightarrow{a},\overrightarrow{b} \rangle = 3\cos 120^{\circ} = - \frac{3}{2}\)；\(\overrightarrow{a}\)在\(\overrightarrow{b}\)上的投影向量为\(\left| \overrightarrow{a} \right|\cos\langle \overrightarrow{a},\overrightarrow{b} \rangle \cdot \frac{\overrightarrow{b}}{\left| \overrightarrow{b} \right|} = - \frac{3}{2} \times \frac{\overrightarrow{b}}{2} = - \frac{3}{4}\overrightarrow{b}\)．
\jans{8}{\ansul{C}}
【分析】根据数量积的定义（零向量情形）、运算律（无结合律）与夹角范围逐一判断．
【详解】解：对A，\(\overrightarrow{a}\)，\(\overrightarrow{b}\)中有零向量时也有\(\overrightarrow{a} \cdot \overrightarrow{b} = 0\)，但零向量与任意向量不谈垂直，A错；对B，数量积是数量，不满足结合律，如\(\overrightarrow{a} \perp \overrightarrow{b}\)且\(\overrightarrow{b} \cdot \overrightarrow{c} \neq 0\)时，左边为零向量而右边不为零向量，B错；对C，\(\left| \overrightarrow{a} \cdot \overrightarrow{b} \right| = \allowbreak\left| \overrightarrow{a} \right| \cdot \allowbreak\left| \overrightarrow{b} \right| \cdot \allowbreak\left| \cos\langle \overrightarrow{a},\overrightarrow{b} \rangle \right| \leq \left| \overrightarrow{a} \right| \cdot \allowbreak\left| \overrightarrow{b} \right|\)，C对；对D，两向量同向共线时夹角为\(0^{\circ}\)，数量积大于0但夹角不是锐角，D错．故选：C．
\jans{9}{\ansul{A}}
【分析】把两式的平方用数量积展开，相反项抵消，结果只与两向量的模有关．
【详解】解：\(\left| \overrightarrow{a} + \overrightarrow{b} \right|^{2} = {\left| \overrightarrow{a} \right|^{2}} + 2\overrightarrow{a} \cdot \overrightarrow{b} + {\left| \overrightarrow{b} \right|^{2}}\)，\(\left| \overrightarrow{a} - \overrightarrow{b} \right|^{2} = {\left| \overrightarrow{a} \right|^{2}} - 2\overrightarrow{a} \cdot \overrightarrow{b} + {\left| \overrightarrow{b} \right|^{2}}\)，两式相加得\(2\left( {\left| \overrightarrow{a} \right|^{2}} + {\left| \overrightarrow{b} \right|^{2}} \right) = 2 \times \left( 4 + 9 \right) = 26\)，结果与夹角无关．故选：A．
\jans{10}{\ansul{证明见解析}}
【分析】连接\(BD\)，用棱向量表示\(\overrightarrow{EH}\)与\(\overrightarrow{FG}\)，证明两向量共线且模不相等．
【详解】证明：连接\(BD\)．因为\(E\)，\(H\)分别是\(AB\)，\(AD\)的中点，所以\(\overrightarrow{EH} = \overrightarrow{EA} + \overrightarrow{AH} = \frac{1}{2}\overrightarrow{BA} + \frac{1}{2}\overrightarrow{AD} = \frac{1}{2}\overrightarrow{BD}\)；因为\(\overrightarrow{CF} = \frac{1}{3}\overrightarrow{CB}\)，\(\overrightarrow{CG} = \frac{1}{3}\overrightarrow{CD}\)，所以\(\overrightarrow{FG} = \overrightarrow{CG} - \overrightarrow{CF} = \frac{1}{3}\left( \overrightarrow{CD} - \overrightarrow{CB} \right) = \frac{1}{3}\overrightarrow{BD}\)．于是\(\overrightarrow{FG} = \frac{2}{3}\overrightarrow{EH}\)，故\(FG \parallel EH\)且\(FG \neq EH\)，四边形\(EFGH\)是梯形．
\jans{11}{\ansul{（1）\(2\sqrt{6}\)，\(2\sqrt{2}\)；（2）\(8\)}}
【分析】设从\(A\)出发的三条棱向量为基底式和链，逐项作数量积展开；展开时注意\(60^{\circ}\)角下两两点积均为正值．
【详解】解：设\(\overrightarrow{AB} = \overrightarrow{a}\)，\(\overrightarrow{AD} = \overrightarrow{b}\)，\(\overrightarrow{AA_{1}} = \overrightarrow{c}\)，则\(\left| \overrightarrow{a} \right| = \allowbreak\left| \overrightarrow{b} \right| = \allowbreak\left| \overrightarrow{c} \right| = 2\)，且\(\overrightarrow{a} \cdot \overrightarrow{b} = \overrightarrow{a} \cdot \overrightarrow{c} = \overrightarrow{b} \cdot \overrightarrow{c} = 2 \times 2\cos 60^{\circ} = 2\)．（1）\(\overrightarrow{AC_{1}} = \overrightarrow{a} + \overrightarrow{b} + \overrightarrow{c}\)，\(\left| \overrightarrow{AC_{1}} \right|^{2} = {\left| \overrightarrow{a} \right|^{2}} + {\left| \overrightarrow{b} \right|^{2}} + {\left| \overrightarrow{c} \right|^{2}} + 2\left( \overrightarrow{a} \cdot \overrightarrow{b} + \overrightarrow{a} \cdot \overrightarrow{c} + \overrightarrow{b} \cdot \overrightarrow{c} \right) = 12 + 2 \times 6 = 24\)，故\(AC_{1} = 2\sqrt{6}\)；\(\overrightarrow{BD_{1}} = \overrightarrow{AD_{1}} - \overrightarrow{AB} = \overrightarrow{b} + \overrightarrow{c} - \overrightarrow{a}\)，\(\left| \overrightarrow{BD_{1}} \right|^{2} = {\left| \overrightarrow{a} \right|^{2}} + {\left| \overrightarrow{b} \right|^{2}} + {\left| \overrightarrow{c} \right|^{2}} + 2\left( \overrightarrow{b} \cdot \overrightarrow{c} - \overrightarrow{a} \cdot \overrightarrow{b} - \overrightarrow{a} \cdot \overrightarrow{c} \right) = 12 + 2 \times \left( - 2 \right) = 8\)，故\(BD_{1} = 2\sqrt{2}\)．（2）\(\overrightarrow{AC_{1}} \cdot \overrightarrow{BD_{1}} = \allowbreak\left( \overrightarrow{a} + \overrightarrow{b} + \overrightarrow{c} \right) \cdot \allowbreak\left( \overrightarrow{b} + \overrightarrow{c} - \overrightarrow{a} \right) = 2\overrightarrow{b} \cdot \overrightarrow{c} + {\left| \overrightarrow{b} \right|^{2}} + {\left| \overrightarrow{c} \right|^{2}} - {\left| \overrightarrow{a} \right|^{2}} = 4 + 4 + 4 - 4 = 8\)．
